\documentclass[UTF8]{ctexart}

\makeatletter
\def\input@path{{../}{../../}}
\makeatother
\usepackage{researchnotes}

\title{补事件与容斥原理}
\author{}
\date{}

\begin{document}
\maketitle

\section{题目与解答}
\label{sec:three-events-problem}

\begin{example}
\label{ex:three-events-original}
设 $A,B,C$ 是样本空间 $\Omega$ 中的三个事件，已知
\[
  P(A)=P(B)=\frac14,\qquad P(C)=\frac25,
  \qquad P(AB)=0,\qquad P(AC)=P(BC)=\frac16.
\]
求 $P(\overline A\,\overline B\,\overline C)$。
\end{example}

\begin{solve}
$AB$ 表示交集 $A\cap B$，即两事件同时发生，并不是概率相乘；$\overline A$ 是 $A$ 的\keyterm{补事件}（complementary event），表示 $A$ 不发生。因此，所求事件是“三个都不发生”，它与“至少有一个发生”的事件 $A\cup B\cup C$ 互为补事件，所以
\begin{equation}
  P(\overline A\,\overline B\,\overline C)=1-P(A\cup B\cup C).
  \label{eq:three-events-complement}
\end{equation}

先由已知条件确定三者交集的概率。三个事件同时发生，一定意味着 $A,B$ 同时发生，即 $ABC=A\cap B\cap C\subseteq AB$。由概率的单调性及条件 $P(AB)=0$，有 $0\le P(ABC)\le P(AB)=0$，所以 $P(ABC)=0$；这里不需要假设事件相互独立。

计算并事件概率要用\keyterm{容斥原理}（inclusion--exclusion principle）：先把三个概率相加，再减去重复计算的两两交集。例如，恰好属于 $A,C$ 的部分被加了两次，减去一次后才保留一次。三者共同部分则被加了三次、减了三次，还要补回一次，使其净计数为 $3-3+1=1$。因此
\begin{equation}
  \begin{aligned}
    P(A\cup B\cup C)
    ={}&P(A)+P(B)+P(C)\\
       &-P(AB)-P(AC)-P(BC)+P(ABC).
  \end{aligned}
  \label{eq:three-events-inclusion-exclusion}
\end{equation}

本题 $P(AB)=P(ABC)=0$，代入式~\eqref{eq:three-events-inclusion-exclusion} 和式~\eqref{eq:three-events-complement}，得
\begin{equation}
  \begin{aligned}
    P(\overline A\,\overline B\,\overline C)
    &=1-\left[\frac14+\frac14+\frac25-\frac16-\frac16\right]\\
    &=1-\frac9{10}+\frac26=\boxed{\frac{13}{30}}.
  \end{aligned}
  \label{eq:three-events-answer}
  \qedhere
\end{equation}
\end{solve}

\end{document}
